\section*{September 4, 2026: Housing controls leave the log--dollar difference largely intact}
\textit{Agent-drafted research entry; not yet reviewed by Jacob.}

Adding housing characteristics reduces the pre-closure fluctuations but leaves
the central pattern intact: post-closure dollar prices rise relative to the
comparison group, while log prices remain near their 2012 relative level.
This comparison uses sales within a quarter mile of the selected 30 closed
and 49 still-open schools on the February 2013 candidate list. It excludes
homes near both groups, near schools designated to receive displaced
students (welcoming schools), or near other schools on the candidate list. All
specifications use the
same transactions and school-site and sale-year fixed effects. Jacob requested
this comparison; both specifications remain exploratory audit results.

The controls account for building and lot size, age, bedrooms, rooms, full
bathrooms, apartment count, property class, residence type, construction
quality, and repair condition. Log building size, log lot size, and age enter
with squared terms. The remaining counts enter linearly and the categories
enter as indicators. These are the corrected Assessor characteristics attached
to each sale year. There are no neighborhood or amenity controls.

\begin{center}
\small
\input{../output/twfe_table.tex}
\end{center}
{\footnotesize Log estimates are reported as $100[\exp(\hat\beta)-1]$, not as
percentage changes in arithmetic mean prices. Dollar estimates use 2022
purchasing power. Pooled DiD compares 2008--2012 with 2014--2018 and excludes
2013. The pre-period test comes from the annual event study relative to 2012.}

The event studies contain 10,921 sales across 29 treated and 48 control sites;
the pooled DiD contains 9,912 sales. Controls do not remove any observations.
The 2010 dip shrinks from $-0.191$ to $-0.120$ log points and from approximately
$-\$31,900$ to $-\$17,700$. The joint pre-period tests no longer reject for
either outcome at conventional levels. This is evidence that observed housing
composition accounts for some of the earlier fluctuation; failure to reject
alone does not establish parallel trends.

The controlled log estimate remains imprecise, with its pooled confidence
interval spanning roughly $-11\%$ to $+25\%$. The pooled dollar interval also
includes zero. The dollar increase survives adjustment for observed housing
quality, so those controls do not explain the difference between outcome
scales. Dollar and log regressions measure different features of the price
distribution: large dollar changes have more influence in the level outcome.
These results do not establish which transactions or parts of the distribution
account for the remaining difference.

\textit{Interpretation.} The controlled estimates are a useful comparison of
prices holding observed housing quality fixed. Condition and improvements may
respond to closure, so this is not automatically the preferred estimate of
the total effect. My recommendation is to retain both specifications while
investigating the remaining distributional and transaction-composition
questions. The evidence does not justify choosing an outcome simply because
its plotted pre-period coefficients look flatter.

\clearpage
\begin{center}
\includegraphics[width=\textwidth]{../input/event_studies.png}
\end{center}
{\footnotesize Top panels omit hedonics; bottom panels include them. Vertical
axes are common across specifications of each outcome. The 2012 coefficient
is normalized to zero; shaded 2013 combines pre- and post-announcement sales.
Intervals are pointwise 95\% confidence intervals clustered by school site.}

\textit{Reproduction.} Run \texttt{make} in
\texttt{tasks/audits/home\_prices\_twfe/code}; build this entry with
\texttt{make} in \texttt{logbook/code}. Tables and the figure link to that
audit's generated outputs through Make. Code revision: \texttt{c200c31} plus
the current uncommitted changes. The local sale and improvement source files
are dated September 3, 2026; the final price sample was rebuilt
September 4 with the approved within-year upper-0.1\% price-per-square-foot
trim. This run did not refresh the raw APIs. The estimator is
\texttt{fixest} 0.12.1 with equal transaction weights and site-clustered
standard errors. Joint pre-period tests use $F(4,76)$. Split-level residence
is redundant with property class 234 and its indicator is dropped.
Independent dummy-variable OLS reproduces all four controlled models'
treatment coefficients within 0.000001, and the uncontrolled estimates match
the prior run.

\clearpage
\section*{Same-day extension: Semiannual event studies}
At Jacob's request, the event studies now split each year into January--June
and July--December. They retain all 10,921 transactions, replace sale-year
with calendar-half-year fixed effects, and use 2012 H2 as the reference.
The annual estimates are unchanged. The difference between outcomes persists:
with hedonics, log estimates fluctuate around zero after closure, while dollar
estimates become more positive in 2016--2018. The joint pre-period tests give
$p=0.967$ for logs and $p=0.572$ for dollars, using $F(9,76)$. Wide intervals
limit what these tests establish. This extension does not explain the
remaining difference between dollar and log outcomes.

\textit{Interpretation clarification.} A flat post-closure log series can still represent an increase from
the years before closure. Equal-weighted half-year averages of the controlled coefficients,
including the normalized reference and excluding 2013, move from $-0.057$ to
$0.001$ in logs and from approximately $-\$5,600$ to $\$50,600$ in dollars.
These are descriptive averages of event coefficients, not new pooled DiD
estimates. Zero denotes the 2012 H2 reference comparison; it does not mean
no change relative to the entire pre-period.

\begin{center}
\includegraphics[width=0.91\textwidth]{../input/event_studies_semiannual.png}
\end{center}
{\footnotesize H1 is January--June; H2 is July--December. Both halves of 2013
are shown: H1 mixes sales before and after the announcement and vote, while
H2 follows the end-of-school-year closure. The nine pre-period coefficients
cover 2008 H1--2012 H1. All other controls, clustering, and sample restrictions
match the annual models. The same audit Make task produces this figure and
exports the semiannual coefficients and model summaries.}
